% Complete proofs for the finite P17-D1 foundation.
\section{Finite preferences and ordinal representation: proofs}
\label{sec:annex-p17-finite-preferences}

Use the definitions of Chapter~\ref{sec:p17-finite-preferences}. The
weak-preference, representation and induced-choice conventions follow
Rubinstein's specified revision
\cite[2015 revision, pp.~5--8, 12--14 and 31]{Rubinstein2015LectureNotes}.
The quotient details, uniqueness on attained values, finite extensions and
worked examples are proved explicitly below.

\begin{annexitem}{annex:p17.preference-structure}{thm:p17-preference-structure}{Indifference and the ordered quotient}
\begin{proof}
\textbf{Relations on alternatives.}\enspace
For any \(x\in X\), completeness applied to \((x,x)\) says
\(x\succeq x\), proving reflexivity. Thus \(x\sim x\).
Symmetry of \(\sim\) follows by exchanging its two defining conjuncts.
If \(x\sim y\) and \(y\sim z\), transitivity gives
\(x\succeq z\) from \(x\succeq y\succeq z\), and
\(z\succeq x\) from \(z\succeq y\succeq x\). Therefore
\(x\sim z\), proving that indifference is an equivalence relation.

The definition of \(x\succ x\) would require both \(x\succeq x\)
and its negation, so \(\succ\) is irreflexive. If \(x\succ y\),
then \(y\succeq x\) is false, whereas \(y\succ x\) would require
it; hence \(\succ\) is asymmetric. If \(x\succ y\) and
\(y\succ z\), weak transitivity yields \(x\succeq z\).
Were \(z\succeq x\) true, it would imply \(z\succeq y\) by
\(x\succeq y\), contradicting \(y\succ z\). Thus \(x\succ z\).
Finally, completeness says at least one weak comparison holds between
\(x,y\). If both hold they are indifferent; if just one holds its
orientation is strict. These three possibilities are exhaustive and
mutually exclusive, including when \(x=y\).

\paragraph{Partition.}
Each \(x\) belongs to \([x]\), so the classes are nonempty and their
union is \(X\). If \([x]\) and \([y]\) meet at \(z\), then
\(z\sim x\) and \(z\sim y\). Symmetry and transitivity imply
\(x\sim y\). For any \(w\in[x]\), transitivity then gives
\(w\sim y\), so \([x]\subseteq[y]\); the reverse inclusion follows
the same way with \(x,y\) exchanged. Thus two classes either coincide
or are disjoint. They form a partition. There are finitely many classes
because every class is \([x]\) for some member of finite \(X\).
For \(X=\varnothing\), the class family is empty and is a partition
of the empty set under these same conditions.

\paragraph{Comparison of classes.}
Suppose \([x]\unrhd[y]\). By definition there are \(a\sim x\)
and \(b\sim y\) with \(a\succeq b\). Then
\(x\succeq a\succeq b\succeq y\), so \(x\succeq y\).
Conversely, if \(x\succeq y\), the members \(x\in[x]\) and
\(y\in[y]\) witness \([x]\unrhd[y]\). This proves the equivalence
for arbitrary representatives.

It transfers reflexivity, completeness and transitivity of \(\succeq\)
to \(\unrhd\). To spell out transitivity, choose representatives
\(a\in A,b\in B,c\in C\). If \(A\unrhd B\) and
\(B\unrhd C\), the equivalence just proved gives
\(a\succeq b\succeq c\), hence \(A\unrhd C\). All classes
are nonempty, so these choices are available. If also
\(A\unrhd B\) and \(B\unrhd A\), representatives \(a,b\)
satisfy \(a\sim b\), so \(A=[a]=[b]=B\). This is antisymmetry.
When there are no classes, all four order conditions hold vacuously.
\end{proof}
\end{annexitem}

\begin{annexitem}{annex:p17.finite-representation}{thm:p17-finite-representation}{Ordering the classes and constructing ranks}
\begin{proof}
\textbf{Necessity.}\enspace
Suppose \(u:X\to\mathbb{R}\) represents a relation \(R\). For every
\(x,y\), at least one of \(u(x)\geq u(y)\) and
\(u(y)\geq u(x)\) holds, proving completeness of \(R\). If
\(xRy\) and \(yRz\), then \(u(x)\geq u(y)\geq u(z)\), so
\(xRz\). Thus \(R\) is transitive. This direction does not require
finiteness.

\paragraph{Least elements in a nonempty finite total order.}
Any nonempty finite subset \(\mathcal A\) of \(\mathcal Q\) has a
least class. Prove this by induction on its positive cardinality. The
singleton case is immediate by reflexivity. For the step, remove one class
\(D\), leaving a nonempty set with a least class \(L\) by induction.
If \(D\unrhd L\), then \(L\) remains below every class in the
full set. Otherwise completeness gives \(L\unrhd D\); since every
remaining class is above \(L\), transitivity places every one above
\(D\). Hence \(D\) is least. If two classes were both least,
each would be above the other, so antisymmetry would make them equal.
This proves existence and uniqueness exactly for nonempty sets.

\paragraph{Enumeration and rank.}
If \(k=|\mathcal Q|>0\), remove its unique least class \(A_0\),
then the unique least class \(A_1\) of what remains, and continue.
At step \(j<k\), exactly \(k-j\) classes remain, so the preceding
paragraph applies. After exactly \(k\) removals none remain. Each
\(A_j\) is below every class removed later. Distinctness and
antisymmetry therefore give
\[
 A_i\unrhd A_j\quad\Longleftrightarrow\quad i\geq j
                                    \qquad(0\leq i,j<k).
\]
In detail, the forward direction cannot hold when \(i<j\), because
then \(A_j\unrhd A_i\) and equality is excluded. The reverse direction
follows from how the least class was selected, or reflexivity when \(i=j\).
Exactly \(A_0,\ldots,A_{j-1}\) are strictly below \(A_j\), so
\(r(A_j)=j\). Thus \(r\) is a bijection onto \(I_k\). The
classes are disjoint and nonempty, so their number is at most \(|X|\).

For \(x,y\in X\), choose their respective class indices \(i,j\).
By the ordered-quotient theorem and the preceding equivalence,
\[
 x\succeq y\quad\Longleftrightarrow\quad
 [x]\unrhd[y]\quad\Longleftrightarrow\quad
 i\geq j\quad\Longleftrightarrow\quad u_0(x)\geq u_0(y).
\]
This proves representation. It also proves correctness and termination of
the class-elimination procedure. A worst alternative in a remaining subset
is exactly a member of its least remaining indifference class: every
alternative in that class is weakly below every remaining alternative, and
an alternative in a higher class is not below the least class. Thus the
source's simultaneous elimination of all worst alternatives removes
\(A_0,A_1,\ldots\) in exactly this order. Its step numbers start at
one, so its constructed values are \(u_0(x)+1\) for nonempty \(X\).

\paragraph{Ties and degenerate cases.}
For any representation \(u\), the two weak comparisons defining
\(x\sim y\) hold exactly when both numerical inequalities hold, which
is equivalent to equality of the values. Likewise, one weak comparison
and failure of the reverse hold exactly when the numerical comparison is
strict. This proves the two displayed equivalences.

If \(X=\varnothing\), then \(X\times X=\varnothing\), so its
only relation is empty and satisfies completeness and transitivity
vacuously. Its only function to \(\mathbb{R}\) is empty and represents that
relation. There are no classes: \(\mathcal Q=I_0=\varnothing\),
and the empty rank function is a bijection between them. The elimination
procedure performs no step and asserts no extremum.

If \(X=\{x\}\), completeness forces its sole pair \((x,x)\)
into the relation; this relation is transitive and any assigned real value
represents it. More generally, on a nonempty all-indifferent set there is
one class, so \(u_0\) is constantly zero. The equality characterization
forces every representation to be constant. Conversely every constant
function gives both weak comparisons for every pair and hence represents
the all-indifferent relation.
\end{proof}
\end{annexitem}

\begin{annexitem}{annex:p17.ordinal-uniqueness}{thm:p17-ordinal-uniqueness}{Scale changes and their exact domain}
\begin{proof}
Suppose first that both \(u,v\) represent the same weak preference.
For \(s\in u(X)\), choose \(x\) with \(u(x)=s\) and prescribe
\(f(s)=v(x)\). This is well defined: if \(u(x)=u(y)\), then
\(x\sim y\), and representation by \(v\) gives \(v(x)=v(y)\).
If \(s<t\) are attained at \(x,y\), respectively, then \(y\succ x\),
so \(v(y)>v(x)\). Thus \(f\) is strictly increasing and hence
injective. Every member of \(v(X)\) is \(v(x)=f(u(x))\) for some
\(x\), so \(f\) is onto \(v(X)\). Its definition proves
\(v=f\circ u\). Any function satisfying that equation must take
\(u(x)\) to \(v(x)\); since these are all the elements of \(u(X)\),
the function is unique on that domain.

Conversely, suppose the stated strictly increasing \(f\) satisfies
\(v=f\circ u\). For any attained \(s,t\), strict increase implies
\[
 s\geq t\quad\Longleftrightarrow\quad f(s)\geq f(t).
\]
The forward implication separates equality from strict inequality. For
the converse, if \(s<t\), strict increase gives the contrary inequality
\(f(s)<f(t)\). Applying this equivalence to \(u(x),u(y)\) proves
that \(v\) represents the same weak comparison as \(u\).

\paragraph{Extending a finite change of scale.}
If the attained domain is empty, its unique map to the empty attained
range extends, for example, to \(F(t)=t\). If it has one element
\(s_0\), with image \(t_0\), use \(F(t)=t-s_0+t_0\).

Otherwise list the attained values increasingly as
\(s_0<\cdots<s_{k-1}\), and put \(t_i=f(s_i)\). Strict increase
gives \(t_0<\cdots<t_{k-1}\). Define a function on all real numbers by
\[
 F(t)\defeq
 \begin{cases}
  t_0+(t-s_0),&t\leq s_0,\\
  t_i+\dfrac{t_{i+1}-t_i}{s_{i+1}-s_i}(t-s_i),
       &s_i\leq t\leq s_{i+1},\ 0\leq i<k-1,\\
  t_{k-1}+(t-s_{k-1}),&t\geq s_{k-1}.
 \end{cases}
\]
At shared endpoints the formulas agree and return the required \(t_i\),
so this is a well-defined extension. Every displayed slope is positive.
For any \(a<b\), split the interval \([a,b]\) at the finitely many
\(s_i\) it contains. Each nonzero subinterval has a positive increase
in \(F\); summing them gives \(F(b)-F(a)>0\). Thus \(F\) is
strictly increasing on \(\mathbb{R}\).

No extension is forced by the finite data outside the attained domain.
For a nonempty attained domain, use slope one or slope two on the ray
strictly below its least point, keeping the same value at that point and
the same strictly increasing construction elsewhere. These give different
strictly increasing extensions with identical values on the attained set;
the same construction works for a singleton. For the empty domain,
\(F(t)=t\) and \(F(t)=t+1\) are already different extensions.

If a strictly increasing \(F:\mathbb{R}\to\mathbb{R}\) is supplied initially,
its restriction to \(u(X)\) preserves and reflects the same numerical
comparisons, so \(F\circ u\) represents the preference. Finally the
two representations in the tied example below are not positive affine
transforms, as its direct calculation proves. Thus affine form is not
a necessary restriction on the scale changes just characterized.
\end{proof}
\end{annexitem}

\begin{annexitem}{annex:p17.maximal-choice}{thm:p17-maximal-choice}{Finite maxima and menu restriction}
\begin{proof}
If \(x\in C_{\succeq}(B)\), then \(x\succeq y\) for every
\(y\in B\), which rules out \(y\succ x\) by the definition of
strict preference. Conversely, suppose \(x\in B\) is strictly
undominated. For any \(y\in B\), if \(x\succeq y\) failed,
completeness would give \(y\succeq x\); together these would say
\(y\succ x\), a contradiction. Hence \(x\succeq y\) for every
\(y\in B\), which proves the first equality. Applying the
representation equivalence separately to each \(y\) proves the second.

For nonempty \(B\), at least one indifference class meets \(B\).
There are finitely many such classes. In the increasing enumeration
\(A_0,\ldots,A_{k-1}\) from the representation proof, let \(j\)
be the largest index of a class meeting \(B\). Every member of
\(B\cap A_j\) is weakly above every alternative in \(B\), by the
class order. If an alternative of \(B\) lies in a lower class, any
member of \(B\cap A_j\) is strictly preferred to it; it therefore
does not belong to the choice set. Thus
\(C_{\succeq}(B)=B\cap A_j\), which is nonempty. This also proves
that all retained alternatives are indifferent. Replacing a representing
function preserves the defining comparisons and hence preserves this set.

For \(B=\varnothing\), the condition that \(x\in B\) leaves no
possible member, so the choice set is empty. For \(B=\{x\}\),
reflexivity supplies its sole required comparison, so the choice set is
\(\{x\}\). More generally, if \(x\succeq y\), reflexivity and
that comparison place \(x\) in \(C_{\succeq}(\{x,y\})\).
Membership in that choice set conversely requires \(x\succeq y\).
This proves recovery, including \(x=y\).

Finally let \(B\subseteq A\) and choose
\(z\in C_{\succeq}(A)\cap B\). Every member of
\(C_{\succeq}(A)\cap B\) is weakly above all elements of \(A\),
and hence all elements of \(B\); this proves its inclusion in
\(C_{\succeq}(B)\). If \(x\in C_{\succeq}(B)\), then
\(x\succeq z\), since \(z\in B\). For any \(y\in A\),
\(z\succeq y\) because \(z\) was chosen from \(A\).
Transitivity gives \(x\succeq y\). Thus
\(x\in C_{\succeq}(A)\cap B\), proving the reverse inclusion.
\end{proof}
\end{annexitem}

\begin{annexitem}{annex:p17.tied-preference-example}{thm:p17-tied-preference-example}{Complete classes and specified menu choices}
\begin{proof}
The relation is defined by the inequalities of the real-valued function
\(t\). Completeness and transitivity of numerical order therefore prove
that it is a weak preference, as in the necessity direction of the
representation theorem. Equal \(t\)-values occur exactly in the three
sets \(\{a,b\},\{c\},\{d,e\}\), so these are all its indifference
classes. Their numerical values show their order. They have respectively
zero, one and two classes below them, proving \(u_0=t\).

The assignment \(0\mapsto0,1\mapsto1,2\mapsto10\) is strictly
increasing on the attained values of \(t\), and its composition with
\(t\) is exactly \(v\). Direct comparison of these three increasing
values gives \(t(x)\geq t(y)\) exactly when \(v(x)\geq v(y)\),
for every pair \(x,y\). Hence \(v\) represents the same relation.

The highest class met by either \(X\) or \(\{a,c,d,e\}\) is
\(\{d,e\}\), and both members are present in each menu. The highest
class met by \(\{a,b,c\}\) is \(\{c\}\), and all members of
\(\{a,b\}\) lie in its single met class. The maximal-choice theorem
therefore yields exactly the four displayed choice sets. Its singleton
and empty-menu clauses give the remaining stated choices.
\end{proof}
\end{annexitem}

\begin{annexitem}{annex:p17.preference-cycle}{thm:p17-preference-cycle}{The complete three-cycle}
\begin{proof}
All three diagonal pairs are listed, so \(R_{\rm cyc}\) is reflexive.
For each pair of distinct alternatives, one of its two orientations is
listed: \((a,b)\), \((b,c)\), or \((c,a)\). Thus the relation is
complete. It contains \((a,b)\) and \((b,c)\), but not \((a,c)\),
so it is not transitive.

Suppose a function \(u\) represented it. The pair \((a,b)\) and
absence of \((b,a)\) would give \(u(a)>u(b)\); the other two
off-diagonal pairs similarly give \(u(b)>u(c)\) and \(u(c)>u(a)\).
Transitivity of strict numerical order would then imply \(u(a)>u(a)\),
which is impossible. Therefore there is no ordinal representation.

The alternative \(a\) fails the comparison against \(c\), \(b\)
fails the comparison against \(a\), and \(c\) fails the comparison
against \(b\). These exhaust the three possible members of
\(C_{R_{\rm cyc}}(X_{\rm cyc})\), so that set is empty. None of
this contradicts the finite-preference choice theorem: its transitivity
hypothesis fails here.
\end{proof}
\end{annexitem}

\begin{annexitem}{annex:p17.ordinal-boundary}{thm:p17-ordinal-boundary}{Gaps and averages are not fixed by an ordinal ordering}
\begin{proof}
The tied-example calculation already proves that \(t\) and \(v\)
represent the same relation. If \(v=\alpha t+\beta\), evaluation at
\(a\), where both functions vanish, gives \(\beta=0\). Evaluation
at \(c\), where both equal one, then gives \(\alpha=1\).
Evaluation at \(d\) would give \(10=2\), a contradiction. Thus
no positive affine transform, or indeed any affine transform, relates
these two functions on all of \(X\).

Substitution of their table entries gives
\(t(c)-t(a)=1-0=1\) and \(t(d)-t(c)=2-1=1\), whereas
\(v(c)-v(a)=1-0=1\) and \(v(d)-v(c)=10-1=9\).
Their endpoint averages are respectively
\((0+2)/2=1=t(c)\) and \((0+10)/2=5>1=v(c)\).
These prove all the numerical assertions. The equality of original
comparisons follows from their representation property, and the
maximal-choice theorem gives equality of their choices on every menu.
The calculation concerns arithmetic on two ordinal representations; it
does not assume or define a preference over lotteries.
\end{proof}
\end{annexitem}
